\begin{table*}[t]
\centering
\caption{证据等级与全文允许措辞。等级表示可证范围，不表示研究价值排序。}
\label{tab:evidence-boundaries}
\footnotesize
\setlength{\tabcolsep}{4pt}
\begin{tabularx}{\textwidth}{@{}p{0.07\textwidth}p{0.20\textwidth}Y Y@{}}
\toprule
等级 & 证据来源 & 允许措辞 & 禁止越界 \\
\midrule
E1 & 冻结提交的控制流、张量运算与配置 & “源码实现/调用/存储/返回”“在该提交中可定位” & 不据此声称有效、稳定、快速或跨模型成立 \\
E2 & 从 E1 公式推出的数学性质与反例 & “一般不保证凸/光滑”“更新秩至多为 $r$”“目标与部署映射不同” & 不把一般性质写成每个样本必然失败，不给出未经证明的收敛率 \\
E3 & 上游说明、代码注释、演示模型与讨论 & “作者报告/注释称/公开演示显示”，并注明未独立复现 & 不写成本文结果，不据单模型演示建立方法排名 \\
E4 & 同行评议文献在其自身协议下的实验 & “该文在指定模型、数据与指标下报告”\cite{piras2026som,arditi2024refusal} & 不跨协议移植数值，不视为矩阵法与基线的同条件对照 \\
E5 & 本文提出但尚未执行的统一协议与假设 & “建议评估/需要检验/可证伪假设”\cite{mazeika2024harmbench} & 不使用“结果表明”“显著提升”“优于”或确定性因果措辞 \\
E6 & 本地 point-v1 冻结配置、journal 标量、日志与机器验收 & “本次 120 个试次完成”“重算得到”“验收失败”\cite{ara2026archive} & 不称独立复现或语义 ASR，不外推通用能力、因果稳定性及 v2 效果 \\
\bottomrule
\end{tabularx}
\end{table*}
